\documentclass{article}
\usepackage{decomp}
\begin{document}
\sourcefile{41598\_2018\_31231\_MOESM1\_ESM.pdf}
\decompnote{The authors' Supplementary Information as supplied by the journal (14 pages, Microsoft Word document printed to PDF, PDF title ``Final\_Supplementary\_Information\_with\_changes\_04-25-2018''). Page numbers (top right of pp.\,2–14) are left out. The Supplementary Movie mentioned in the main text is a separate file and is not part of this PDF.}

\pdfpage{1}
\noindent 04-25-18

\begin{center}
\textbf{\large Supplementary Information}

\bigskip
\textbf{\large Biological serial block face scanning electron microscopy at improved z-resolution based on Monte Carlo model}
\end{center}

\noindent Q. He\textsuperscript{1}, M. Hsueh\textsuperscript{2}, G. Zhang\textsuperscript{1}, D.C. Joy\textsuperscript{3, 4}, R.D. Leapman\textsuperscript{1*}

\noindent\textsuperscript{1}National Institute of Biomedical Imaging and Bioengineering, NIH, Bethesda, MD, 20892, USA.\\
\textsuperscript{2}Gatan, Inc., 5794 W. Las Positas Blvd, Pleasanton, CA, 94588, USA\\
\textsuperscript{3}Department of Material Science and Engineering, University of Tennessee, Knoxville, TN, 37996, USA\\
\textsuperscript{4}Center for Nanophase Material Sciences, Oak Ridge National Laboratory, Oak Ridge, TN, 37831, USA

\bigskip
\noindent\textsuperscript{*}Corresponding author:\\
Richard D. Leapman\\
National Institute of Biomedical Imaging and Bioengineering,\\
Building 13, Room 3N17,\\
13 South Drive,\\
Bethesda, MD 20892, USA\\
Tel: 301-496-2599\\
Email: leapmanr@mail.nih.gov

\pdfpage{2}
\begin{center}
\includegraphics[width=0.5\linewidth,height=0.4\textheight,keepaspectratio]{figures/sfig1.png}
\end{center}
\noindent\textbf{Supplementary Figure 1 |} Representation of image stack used for sub-slice reconstruction. White rectangles indicate a continuous column structure that penetrates the sample. The slight shifts of columns in different cuts are due to residual charge-induced electrostatic beam deflection. Sub-slice reconstruction requires that image alignment be accurate to within a single pixel.
\alttext{Schematic titled ``image stack'': a vertical stack of eight thin slabs drawn in perspective, in pairs of shades: dark blue and light blue (labelled ``cut 1'' at the left), red and crimson (``cut 2''), green and light green (next to a vertical ellipsis), purple and light purple (``cut n''). The slabs are labelled alternately at the right ``E\textsubscript{low}'' (upper, darker slab of each pair) and ``E\textsubscript{high}'' (lower, lighter slab). Each slab shows a small white-outlined rectangle near its middle (the column structure of the caption): a short, wide one on the E\textsubscript{low} slabs and a taller, narrow one on the E\textsubscript{high} slabs; their horizontal positions differ slightly from slab to slab.}

\pdfpage{3}
\begin{center}
\includegraphics[width=0.95\linewidth,height=0.75\textheight,keepaspectratio]{figures/sfig2.png}
\end{center}
\pdfpage{4}
\noindent\textbf{Supplementary Figure 2 |} Backscattered images acquired from six successive 25-nm cuts of a liver sample that had been stained with heavy metals and embedded in epoxy resin. The block-face images are shown after fine alignment to correct for drift. (a), (c) BSE images acquired at beam energies of \emph{E\textsubscript{low}} = 1.0 keV and \emph{E\textsubscript{high}} = 1.4 keV; (b), (d) pre-processed BSE images in (a) and (c) after setting negative pixels intensities to zero. (e) (f) and (g) pre-processed BSE images and difference images of cut 1, 3 and 5 from (b) and (d) after normalizing the background. To display differences between BSE images acquired at dual landing energies, pairs of images in Figs. b1 / d1, b3 / d3, and b5 / d5 are shown at higher magnification in Figs. eI / eII, fI / fII, and gI / gII, respectively, and their differences are displayed in Figs. eIII, fIII, and gIII. White circles indicate typical differences between features, which appear bright or dark depending on whether they are located in the top or bottom subslice. Red triangles indicate homogenous region used for scaling image intensities in the \emph{E\textsubscript{low}} and \emph{E\textsubscript{high}} image stacks. Scale bar, 500 nm.
\alttext{Grey-scale BSE images of liver. Rows a–d: four rows of six small images each, numbered ``1''–``6'' in large white digits at the bottom (the cut number), each with a white scale bar at the upper right. All show the same field: several parallel, wavy bright membrane lines running from upper left to lower right, a bright irregular patch left of centre, and a round, finely granular region at the lower left. Rows a and c (raw images at the two energies) have a mid-grey background; rows b and d (pre-processed) have a black background with the bright lines standing out. A small red open triangle sits at the lower left of image a1 and of image c1, on the round granular region. Rows e–g: three rows of three larger images labelled ``I'', ``II'', ``III'' in white at the bottom and ``e'', ``f'', ``g'' at the upper left of the first image; each has a white scale bar at the upper right. I and II are enlarged pre-processed images (black background, bright speckled membrane lines); III is a mid-grey noisy difference image in which the membranes are barely visible. Two to four small white circles mark the same positions in I, II and III of each row, around small bright spots in I or II.}

\pdfpage{5}
\begin{center}
\includegraphics[width=0.95\linewidth,height=0.6\textheight,keepaspectratio]{figures/sfig3.png}
\end{center}
\noindent\textbf{Supplementary Figure 3 |} Four pairs of backscattered images at (a) \emph{E\textsubscript{low}} = 1.0 keV and (b) \emph{E\textsubscript{high}} = 1.4 keV from successive 25-nm cuts of a heavy metal-stained liver block, together with calculated structure of (c) top and (d) bottom sub-slices, each of thickness 12.5 nm. The numbers indicated in white represent the cutting order and correspond to the numbers in Supplementary Fig. 2. The numbered arrows (red and blue) denote selected fine-scale features that change from sub-slice to sub-slice. Scale bar, 500 nm.
\alttext{Sixteen grey-scale BSE images of the same liver field (wavy bright membrane lines on a dark background, a round granular region at the lower left), in a 4 × 4 grid. Each image carries a letter at the upper left (``a'', ``b'', ``c'', ``d''), a large white cut number at the bottom (``5'', ``6'', ``7'', ``8'') and a short white scale bar at the lower left. Left half, top: cut 5 (a, b above; c, d below), with a red arrow labelled ``1'' pointing down-right at a feature right of centre. Right half, top: cut 6 (a, b; c, d), with a blue arrow labelled ``2'' pointing up-right near the upper right. Left half, bottom: cut 7, with a blue arrow ``4'' pointing up-right and a red arrow ``3'' pointing down-right below it. Right half, bottom: cut 8, with a red arrow ``5'' pointing right near the upper right. In each block of four, a and b (the measured images at the two energies) look smoother; c and d (the calculated top and bottom sub-slices) are noisier and more contrasted, and the feature at the arrow looks different between c and d.}

\pdfpage{6}
\begin{center}
\includegraphics[width=0.85\linewidth,height=0.45\textheight,keepaspectratio]{figures/sfig4.png}
\end{center}
\noindent\textbf{Supplementary Figure 4} | Energy-dispersive x-ray spectrum from thin section of tissue block used for SBF-SEM imaging, recorded at 120 keV beam energy using an FEI Tecnai TEM equipped with an Oxford Instruments X-Max SDD and Inca microanalysis system. X-ray spectra generated by the DTSA-II software were used as reference spectra, which were fitted to the tissue spectrum to quantitate the relative concentrations of the heavy atoms: Os, Pb, and U.
\alttext{Line plot (low-resolution raster image). x axis ``X-ray Energy (keV)'', linear, 8 to 20 (labels every 2); y axis ``Counts'', linear, 0 to about 700 (labels 0, 200, 400, 600). Legend at the upper right: ``Sectioned biological sample block'' (black), ``Copper standard'' (blue), ``Lead standard'' (purple), ``Osmium standard'' (green), ``Uranium standard'' (red). The black spectrum has a noisy baseline of about 10–30 counts and peaks labelled in the plot: at about 8.9 keV, labelled ``Cu K–M2 / Os L3–M4'', the tallest peak (about 590 counts), with the blue (about 280) and green (about 255) standards peaking at the same energy below it; a double peak labelled ``Os L2–M4'' (arrow to the shoulder at about 10.35 keV, about 120) and ``Os L1–M3 / Pb L3–M5'' (arrow to the peak at about 10.55 keV, about 215), where the green standard peaks at about 10.35 keV (about 100) and the purple at about 10.55 keV (about 145); ``Pb L2–M4'' at about 12.6 keV (black about 110, purple about 80); ``U L3–M5'' at about 13.6 keV (black about 55, red about 45). Smaller black bumps (about 25–30 counts) near 12.1, 14.8, 16.4 and 17.2 keV, matched by small peaks of the green (12.1 keV), purple (14.8 keV) and red (16.4 and 17.2 keV) standards. The blue Copper standard rises steeply at the left edge (about 8.2 keV) to above the top of the plot. The other standards lie near 0 away from their peaks. Values from sfig4.csv.}
\figuredata{(single panel)}{figures/sfig4.csv}{figures/sfig4.py}{digitized from the embedded raster image, about 84 ppi (tier 3)}

\bigskip
\noindent The x-ray spectrum shows that the sample contains uranium, lead and osmium, whereas the high copper peaks are due to x-ray generation in the copper grid used to mount the thin sections, and these are ignored in the analysis. There is strong overlap between the osmium L lines with copper K lines at \textasciitilde{} 8.9 keV, and between the lead L lines and the osmium L lines, which complicates quantification of these elements. We therefore used DTSA-II\textsuperscript{1} to fit the spectrum, as described in Methods. The relative ratio of heavy atoms was: lead 0.0561±0.0009, osmium 0.0507±0.0003, and uranium 0.0193±0.0007. Thus, lead is the dominant stain in this specimen.

\pdfpage{7}
\begin{center}
\includegraphics[width=0.95\linewidth,height=0.6\textheight,keepaspectratio]{figures/sfig5.png}
\end{center}
\noindent\textbf{Supplementary Figure 5 |Effect of stain composition on the linearity of the backscattered signal}. \textbf{(a)} Three-dimensional geometrical model for 50 nm x 50 nm x 12.5 nm lead stained cuboids located at different depths in a 800 nm by 800 nm by 800 nm epoxy block and its view of x-z plane, y-z plane and x-y plane. Centers of cuboids are located 6.25 nm, 18.75 nm, 12.5 nm, respectively, from the top surface of the block from left to right. Dimension in z is not drawn to scale. \textbf{(b)} Intensity vs. stain composition profiles for cuboids shown in (a), simulation data \pdfpage{8}were calculated at 1.0 keV with a stain composition from 1\%-9\%. \textbf{(c)} Intensity vs. stain composition profiles for cuboids shown in (a); simulation data were calculated at 1.4 keV with a stain composition from 1\%-9\%. \textbf{(d)} Intensity vs. stain composition profiles for cuboids shown in (a); simulation data were calculated at 1.0 keV with a stain composition from 10\%-100\%. \textbf{(e)} Intensity vs. stain composition profiles for cuboids shown in (a); simulation data were calculated at 1.4 keV with a stain composition from 10\%-100\%. Curves were fitted though the data using a second order polynomial regression. The red curve is based on the intensity of cuboid 3 in (a), and the blue curve is based on the sum of the backscattered signals from cuboid 1 and cuboid 2. The curves in (a) and (c) show that the backscattered signal depends linearly on stain content for concentrations below 5 atomic percent, whereas the curves in (c) and (d) reveal a strong non-linear dependence on stain content for concentrations above 10 atomic percent.
\decompnote{The caption's last sentence refers to ``(a) and (c)'' and ``(c) and (d)''; by the panel definitions above, the 1–9\,\% curves are in (b) and (c) and the 10–100\,\% curves in (d) and (e). The caption also calls the ``sum'' curve ``blue''; in the figure the sum markers are blue diamonds and the ``Poly. (sum)'' trendline is purple. Kept as printed.}
\alttext{(a) Schematic: a wire-frame cube (dashed hidden edges) with three cuboids just below its top face, numbered ``1'', ``2'', ``3'' from left to right: 1 and 2 blue (2 lower than 1), 3 red and taller, spanning both depths. Short arrows of the same colour point at each cuboid from above and below. Colour key at the left: ``0-12.5 nm'' (blue), ``12.5-25 nm'' (blue), ``0-25 nm'' (red). Axis arrows at the lower right: ``z'' up, ``y'' into the page, ``x'' to the left. Three framed views at the right: x-z (``z'' up, ``x'' right): blue square at the upper left, a second blue square lower, a taller red rectangle at the right; y-z (``z'' up, ``y'' right): two blue squares stacked; x-y (``y'' up, ``x'' right): blue, blue, red squares in a row. (b)–(e) Four scatter charts, each titled at the top: (b) ``low stain level 1.0 keV'', (c) ``low stain level 1.4 keV'', (d) ``high stain level 1.0 keV'', (e) ``high stain level 1.4 keV''. x axis ``stain level (atomic percent)'': 1\% to 9\% (labels 1\%, 3\%, 5\%, 7\%, 9\%) in b and c, 10\% to 100\% (labels 10\%, 40\%, 70\%, 100\%) in d and e; y axis ``counts (a.u.)'': 0–1600 (b), 0–2000 (c, d), 0–3000 (e), with horizontal gridlines. Legend at the right of each: ``0-25 nm'' (red triangles), ``sum'' (blue diamonds), ``Poly. (0-25 nm)'' (red line), ``Poly. (sum)'' (purple line). In b, both series rise nearly linearly from about 220–240 at 1\% to about 1210–1350 at 9\%, the red triangles slightly above the blue diamonds. In c, both rise from about 280 at 1\% to about 1470–1580 at 8–9\%, almost on top of each other. In d, both rise steeply and level off: red from about 700 at 10\% to about 1550 at 100\%, blue from about 620 to about 1780, the two crossing between 20\% and 30\%. In e, red rises from about 750 to about 1750 and flattens from about 60\% on, blue from about 800 to about 2630, clearly above red from 20\% on. Values from sfig5b–e.csv.}
\figuredata{b}{figures/sfig5b.csv}{figures/sfig5b.py}{vector paths and placed marker images in the PDF, exact up to marker placement (tier 2)}
\figuredata{c}{figures/sfig5c.csv}{figures/sfig5c.py}{vector paths and placed marker images in the PDF, exact up to marker placement (tier 2)}
\figuredata{d}{figures/sfig5d.csv}{figures/sfig5d.py}{vector paths and placed marker images in the PDF, exact up to marker placement (tier 2)}
\figuredata{e}{figures/sfig5e.csv}{figures/sfig5e.py}{vector paths and placed marker images in the PDF, exact up to marker placement (tier 2)}

\begin{center}
\small
\begin{tabular}{p{0.62\linewidth}p{0.25\linewidth}}
\toprule
\multicolumn{1}{c}{\textbf{Parameter}} & \\
\midrule
\textbf{Number of incident electrons / pixel} & 10, 000\\
\textbf{Electron beam diameter} & 12.5 nm\\
\textbf{Pixel size} & 12.5 nm\\
\textbf{Simulation noise} & Shot noise\\
\textbf{Backscattered detector efficiency (}$\boldsymbol{\epsilon_{det}}$\textbf{)} & 1.0\\
\textbf{Low primary energy for stained cuboids in epoxy (25 nm/cut)} & 1.0 keV\\
\textbf{High primary energy for stained cuboids in epoxy (25 nm/cut)} & 1.4 keV\\
\textbf{Low primary energy for stained cuboids in epoxy (50 nm/cut)} & 1.4 keV\\
\textbf{High primary energy for stained cuboids in epoxy (50 nm/cut)} & 2.2 keV\\
\textbf{Density of epoxy embedding resin} & 1.22 g/cm\textsuperscript{3}\\
\textbf{Composition of epoxy resin} & C\textsubscript{36}H\textsubscript{52}O\textsubscript{12}\\
\textbf{Heavy atom stain concentration (Pb)} & 3 atomic \%\\
\textbf{Pb stain density in cuboids} & 3.24 atoms/nm\textsuperscript{3}\\
\bottomrule
\end{tabular}
\end{center}
\noindent\textbf{Supplementary Table 1}. Parameters used in Monte Carlo simulation.
\decompnote{In the PDF the table rows are shaded alternately grey and white. Also in tables/stable1.csv. The text layer of the PDF garbles the symbol $\epsilon_{det}$ (``!!"\#''); it is transcribed from the page image.}

\pdfpage{9}
\subsection*{Supplementary Notes}

\noindent\textbf{Supplementary Note 1:} To test the sensitivity of the sub-slice reconstruction technique, a model structure was generated containing two cuboids of size 50 nm × 50 nm × 25 nm with a lead concentration of 3\% in a pure epoxy resin matrix, with one cuboid extending from the block surface to a depth of 25 nm, and the other cuboid from 25 nm below the block surface to 50 nm below the surface. Since this simulated dataset corresponded to a cutting increment of 50 nm, an electron fluence of 20 e/nm\textsuperscript{2} was considered (i.e., somewhat higher than the maximum electron fluence of 15 e/nm\textsuperscript{2} for cutting at increments of 25 nm). Simulated image intensities generated by the cuboids for beam energies of 1.4 keV and 2.2 keV were first processed to subtract the background, which gave the coefficients of the 2×2 matrix:
\[
\mathbf{A}\ = \begin{pmatrix} 0.757{\times}10^{-3} & 0.014{\times}10^{-3}\\ 0.598{\times}10^{-3} & 0.400{\times}10^{-3} \end{pmatrix}
\]
from which the inverse matrix can be calculated:
\[
\mathbf{A}^{-1} = \begin{pmatrix} 1.359{\times}10^{3} & -0.048{\times}10^{3}\\ -2.031{\times}10^{3} & 2.571{\times}10^{3} \end{pmatrix}
\]
\decompnote{The two matrices are set in Cambria Math in the PDF and its text layer garbles them (``!'' for the exponents); they are transcribed from the page image.}

\pdfpage{10}
\bigskip
\noindent\textbf{Supplementary Note 2: DM script used to take the dual energy dataset.}

\decompnote{The script is reproduced line by line from the PDF (pp.\,10–14) in a verbatim block, including its comments and inconsistencies (e.g.\ comments that do not match the energies set in the code).}
\begin{verbatim}
// $BACKGROUND$
//define total cuts
number ncut = 80
number SCOPE_DELAY = 8
//number delta_x = 42.033
//number delta_y = 14.856
//number delta_beamshiftx= 13.802
//number delta_beamshifty = -20.658
number NewObstigx = 10.821
number NewObstigy = -1.683
number NewFocus =6008
number NewMag =8950
number Mag =8750
number NewStageX =-394.79
number NewStageY =-205.6
//number delta_focus = -882
//number delta_mag = 1003
//cutting and imaging loop start
for (number n = 0; n < ncut; n++)
{
EMWaitUntilReady( )
EMUpdateCalibrationState( )
Result( "\n\n\n" )
Result( " **** Basic Microscope Parameters **** \n" )
if ( !EMIsReady( ) )
{
Result( "--waiting for microscope to be ready-- \n" )
EMWaitUntilReady( )
}
Result( "--microscope is ready-- \n\n" )
EMUpdateCalibrationState( )
Result( "Microscope: " + EMGetMicroscopeName( ) + "\n" )
//read microscope high tension
if ( EMCanGetHighTension( ) )
{
Result( "High Tension: " + EMGetHighTension( ) + " Volt \n" )
}
else
{
Result( "High Tension: CAN NOT BE READ \n" )
}
\end{verbatim}
\pdfpage{11}
\begin{verbatim}
Result( "Operation mode: " + EMGetOperationMode( ) + "\n" )
//read microscope magnification
if ( EMCanGetMagnification( ) )
{
if ( EMGetOperationMode( ) != "DIFFRACTION \n" )
{
Result( " -Magnification: " + EMGetMagnification( ) + "\n")
}
}
else
{
Result( " -Magnification: CAN NOT BE READ \n" )
}
// read focus and stigmation
number focus // focus
number ObStigX, ObStigY // objective stigmator
//number Brightness, Contrast
//void PrintAllValues( number calibrated )
Result( "\n\n" )
Result( "--- RAW VALUES --- \n" )
focus = EMGetFocus( )
EMGetObjectiveStigmation( ObStigX, ObStigY )
//EMGetBrightness(Brightness)
Result( "Focus : " + focus + "\n" )  //focus is 1000 times bigger than the real value
Result( "Obj. Stig : " + ObStigX + "/"+ ObStigY + "\n" )
EMWaitUntilReady( )
// read stage position (x,y,z)
Result( "\n\n" )
Result( "Current Stage Positions: \n" )
number stageX, stageY, stageZ
stageX = EMGetStageX()
stageY = EMGetStageY()
stageZ = EMGetStageZ()
Result( "3viewStageX: " + stageX + "\n" )
Result( "3viewStageY: " + stageY + "\n" )
Result( "3viewStageZ: " + stageZ + "\n" )
EMWaitUntilReady( )
//read current beam shift value
number x_beam,y_beam
EMGetBeamShift(x_beam,y_beam)
Result( "current beam shift " +x_beam +"and y" + y_beam + "\n" )
//set beam energy value to 1.4 keV
number BeamEnergy1 = 1400  //notice the unit here is volt
EMSetBeamEnergy(BeamEnergy1)
EMWaitUntilReady( )
//acquire image at 1.4 keV
\end{verbatim}
\pdfpage{12}
\begin{verbatim}
EMSetBeamBlanked(1)  //1 means beam blank
Sleep(SCOPE_DELAY)
EMSetBeamBlanked(0) //0 means beam unblank
//image img+n
image img14 :=integerimage("Imaging at" +"1.4keV"+n, 2,0,2000,2000)
EMSetBeamBlanked(1)  //1 means beam blank
Sleep(SCOPE_DELAY)
EMSetBeamBlanked(0) //0 means beam unblank
DSAcquireData(img14,1,4,0,0)
showImage(img14)
Result( "acquire image at 1.0kev DONE \n" )
EMWaitUntilReady( )
//save image at 1.4 keV
//Image fImage = GetFrontImage()
SaveAsGatan(img14, "E:\\11142017\\"+ n + "lowKeV")
Result( "IMAGE SAVE 1.4KEV done \n" )
EMWaitUntilReady( )
Result( "debug1" )
// move stage position (x,y,z)
Result( "\n\n" )
Result( "moved Stage Positions: \n" )
//number stageX1 = stageX + delta_x
//number stageY1 = stageY + delta_y
number stageZ1 = stageZ
EMSetStageX(NewStageX)
EMSetStageY(NewStageY)
EMSetStageZ(stageZ1)
Result( "3viewStageX1: " + NewStageX + "\n" )
Result( "3viewStageY1: " + NewStageY+ "\n" )
Result( "3viewStageZ1: " + stageZ1 + "\n" )
EMWaitUntilReady( )
//set beam energy value to 1.4 keV
number BeamEnergy3 = 2200
EMSetBeamEnergy(BeamEnergy3)
EMWaitUntilReady( )
 //set new mag
//number magindex = EMGetMagIndex()
//EmSetMagIndex (magindex + 1)
EMWaitUntilReady()
//number Mag = EMGetMagnification( )
//number Newmag = Mag + delta_mag
EMSetMagnification(NewMag)
EMWaitUntilReady( )
number Newmagv = EMGetMagnification( )
Result( "changed magnification: " + newMagv + "\n" )
//set new focus
\end{verbatim}
\pdfpage{13}
\begin{verbatim}
//number oldfocus, newfocus
//oldfocus = EMGetFocus()
//newfocus = oldfocus + delta_focus
EMSetFocus(NewFocus)
EMWaitUntilReady( )
//AFS_Run()
//set new objstigx_y
//number newObStigX, newObStigY
number NObStigX, NObstigY
//newObStigX = ObStigX + delta_stigx
//newObStigY = ObStigY + delta_stigy
EMSetObjectiveStigmation( newObStigX, newObStigY )
EMGetObjectiveStigmation( NObStigX, NObStigY )
Result( "changed stigmation: " + NobStigX + "/" + NObStigY +"\n" )
EMSetBeamBlanked(1)  //1 means beam blank
Sleep(SCOPE_DELAY)
EMSetBeamBlanked(0) //0 means beam unblank
//EMSetObjectiveStigmation( newObStigX, newObStigY )
//read current beam shift value
number x_beam1,y_beam1
EMGetBeamShift(x_beam1,y_beam1)
Result( "current beam shift2 " +x_beam1 +"and y" + y_beam1 + "\n" )
//acquire image at 1.4 keV
image img34 :=integerimage("Imaging at" +"2.2keV"+n, 2,0,2000,2000)
EMSetBeamBlanked(1)  //1 means beam blank
Sleep(SCOPE_DELAY)
EMSetBeamBlanked(0) //0 means beam unblank
DSAcquireData(img34,1,2,0,0)
showImage(img34)
//save image at 3.4 keV
//Image fImage1 = GetFrontImage()
SaveAsGatan(img34, "E:\\11142017\\" + n + "highKeV")
//raise the sample up
number deltaZ = 0.05 //um
number stageZnew = stageZ + deltaZ
EMSetStageZ (stageZnew)
Result( "3viewStageZmoved: " + stageZnew+ "\n" )
EMSetBeamBlanked(1)  //1 means beam blank
Sleep(SCOPE_DELAY)
EMSetBeamBlanked(0) //0 means beam unblank
//cut the sample and blank the beam in this process
Number StrokeUp = 64 + 32 + 16 + 8 + 4 + 2*0 + 1*0
Number StrokeDown = 64 + 32 + 16 + 8 + 4 + 2*1 + 1*1
EMSetBeamBlanked(1)  //1 means beam blank
MicrotomeStage_Near()
DSSetDigitalOutput(StrokeUp)
\end{verbatim}
\pdfpage{14}
\begin{verbatim}
MicrotomeStage_cut(1)
MicrotomeStage_Retract(1)
DSSetDigitalOutput(StrokeDown)
MicrotomeStage_Clear()
Sleep(SCOPE_DELAY)
EMSetBeamBlanked(0) //0 means beam unblank
Result( "cut process skipped " + "\n" )
//set beam energy value back to 1.0 keV
number BeamEnergy4 = 1400
EMSetBeamEnergy(BeamEnergy4)
// move stage position back (x,y,z)
Result( "\n\n" )
Result( "moved Stage Positions: \n" )
EMSetStageX(stageX)
EMSetStageY(stageY)
EMSetStageZ(stageZnew)
Result( "3viewStageX_old_value: " + stageX + "\n" )
Result( "3viewStageY_old_value: " + stageY + "\n" )
Result( "3viewStageZ_old_value: " + stageZnew + "\n" )
//set Magnification back
EMSetMagnification(Mag)
//set ObstigX and Y back
EMSetObjectiveStigmation(ObStigX, ObStigY )
//set focus back
EMSetFocus(focus)
//AFS_Run()
}
//turn off Beam energy when done
EMSetBeamBlanked(1)  //1 means beam blank
EMSetHighTensionOnOff(0)  // '1/0' = 'on/off'
\end{verbatim}


\subsection*{Supplementary References}

\noindent 1.\quad Ritchie N. W. M. Spectrum simulation in DTSA-II. \emph{Microsc Microanal} \textbf{15}, 14 (2009).
\decompnote{The main article cites the same paper as ``\emph{Microsc. Microanal.} \textbf{15}, 454–468 (2009)'' (ref.\ 43); the supplement gives ``15, 14''.}

\end{document}
